\documentclass[a4paper]{article}
% Español (México) con XeLaTeX: fontspec para fuentes Unicode, polyglossia
% para la división silábica y los nombres del documento en la variante mexicana.
\usepackage{fontspec}
\usepackage{polyglossia}
\setdefaultlanguage[variant=mexican]{spanish}
% Los nombres chinos van como «pinyin (original)»: xeCJK da a los caracteres
% Han su propia fuente; sin ella XeLaTeX los omite con un aviso.
% FandolSong no cubre algunos caracteres de nombres (U+73FA); WenQuanYi Zen Hei sí.
\usepackage[AutoFallBack=true]{xeCJK}
\setCJKmainfont{FandolSong-Regular.otf}
\setCJKfallbackfamilyfont{\CJKrmdefault}{WenQuanYi Zen Hei}
% polyglossia no traduce los nombres de listings.
\AtBeginDocument{\providecommand{\lstlistingname}{}\renewcommand{\lstlistingname}{Listado}%
  \providecommand{\lstlistlistingname}{}\renewcommand{\lstlistlistingname}{Índice de listados}}
\usepackage{amsmath, amssymb}
\usepackage{graphicx}
\usepackage[margin=2.35cm]{geometry}
\usepackage[most]{tcolorbox}
\usepackage{booktabs}
\usepackage{tabularx}
\usepackage{longtable}
\usepackage{array}
\usepackage{float}
\usepackage{xcolor}
\usepackage{hyperref}
\usepackage{fancyhdr}

\hypersetup{colorlinks=true, linkcolor=blue!55!black, urlcolor=blue!60!black}
\graphicspath{{./}{figures/}}
\setlength{\parskip}{0.55em}
\setlength{\parindent}{2em}
\setlength{\emergencystretch}{3em}
\renewcommand{\arraystretch}{1.18}

\newtcolorbox{knowledgebox}[1]{enhanced, colback=blue!5!white, colframe=blue!70!black, colbacktitle=blue!70!black, coltitle=white, fonttitle=\bfseries, title={#1}, sharp corners, breakable}
\newtcolorbox{importantbox}[1]{enhanced, colback=yellow!10!white, colframe=yellow!75!black, colbacktitle=yellow!75!black, coltitle=black, fonttitle=\bfseries, title={#1}, sharp corners, breakable}
\newtcolorbox{warningbox}[1]{enhanced, colback=red!5!white, colframe=red!70!black, colbacktitle=red!70!black, coltitle=white, fonttitle=\bfseries, title={#1}, sharp corners, breakable}

\pagestyle{fancy}
\fancyhf{}
\fancyfoot[C]{\thepage}
\fancyfoot[R]{\footnotesize\color{gray}\href{https://github.com/hqhq1025/ai-course-notes}{AI Course Notes}}
\renewcommand{\headrulewidth}{0pt}
\renewcommand{\footrulewidth}{0pt}

\newcommand{\videourl}{https://www.youtube.com/watch?v=vG1RBqn1sG4}
\newcommand{\repourl}{https://github.com/hqhq1025/ai-course-notes}

\begin{document}
\begin{titlepage}
\centering
{\Large Notas de entrevista técnica\par}
\vspace{1.0cm}
{\huge\bfseries Entrevista con Luo Fuli (罗福莉): el paradigma Agent, el posentrenamiento y la reestructuración organizacional de los laboratorios de modelos\par}
\vspace{0.6cm}
{\Large Zhang Xiaojun conversa con Luo Fuli\par}
\vspace{0.4cm}
{\large Elaborado a partir del video público del Zhang Xiaojun Podcast, la transcripción hecha en GPU y diagramas conceptuales propios\par}
\vspace{0.3cm}
{\large 2026-04-24\par}
\vspace{0.9cm}
\includegraphics[width=0.88\textwidth,height=0.42\textheight,keepaspectratio]{cover.jpg}\par
\vfill
\begin{tcolorbox}[width=0.92\textwidth, colback=black!2!white, colframe=black!55, sharp corners]
\textbf{Título del video}: 138. Entrevista de 3.5 horas con Luo Fuli: ¡el paradigma de la IA ya cambió radicalmente! OpenClaw, el paradigma Agent depende mucho del posentrenamiento, la asignación de GPU y la igualdad dentro de la organización\par
\textbf{Canal del video}: Zhang Xiaojun Podcast\par
\textbf{Duración del video}: 03:36:36\par
\textbf{Enlace del video}: \href{\videourl}{\nolinkurl{\videourl}}\par
\textbf{Notas sobre el material}: YouTube no ofrece subtítulos; el texto se elaboró a partir de la transcripción con faster-whisper large-v3 en CUDA, los capítulos del video, la portada y figuras didácticas propias. Las tomas fijas de la entrevista no se reutilizan en el texto.\par
\textbf{Página del proyecto}: \href{\repourl}{\nolinkurl{\repourl}}\par
\end{tcolorbox}
\end{titlepage}

\tableofcontents
\newpage
\section{Introducción: por qué importa esta entrevista}

Esta entrevista reúne en una misma conversación un conjunto de palabras clave de la industria de los grandes modelos a inicios de 2026: OpenClaw, Agent, postentrenamiento, RL scaling, MiMo-V2, orquestación de múltiples modelos, modelo base de 1T, asignación de GPU, igualdad organizacional, el entorno importa más que la experiencia. La invitada, Luo Fuli (罗福莉), trabajó en Alibaba DAMO Academy (阿里达摩院) y en DeepSeek, y actualmente dirige el equipo de grandes modelos de Xiaomi. El valor de la entrevista no está solo en «qué tan bien se desempeña cierto modelo», sino en que muestra cómo un AI Lab percibe un cambio de paradigma y reorganiza su forma de investigar.

\begin{importantbox}{Tesis central de este episodio}
El juicio de Luo Fuli es que la guerra de los grandes modelos entra en su segundo acto: de la era del Chat, dominada por el pre-train, a la era del Agent, dominada por el post-train. En adelante, la competencia no dependerá solo del modelo base, sino también del framework de Agent, la RL infra, los entornos de tareas, el costo y la velocidad, la agilidad organizacional y la asignación de capacidad de cómputo.
\end{importantbox}

\begin{knowledgebox}{Nota sobre la estrategia visual}
Este video es una entrevista con plano fijo, sin slides didácticas, pizarrón ni demostraciones de producto. Conforme al estándar de pódcast de este repositorio, el cuerpo del texto no repite fotogramas de las personas; la portada sirve para identificar la fuente, y el contenido técnico se presenta mediante diagramas conceptuales y tablas.
\end{knowledgebox}

\subsection{Resumen de la sección}

Este episodio debe leerse como material oral de un «manual de operación de un AI Lab en la era del Agent»: no trata solo de modelos, sino de cómo cambian en conjunto los modelos, los frameworks, el postentrenamiento, la capacidad de cómputo y la organización.
\section{De Chat a Agent: un cambio de paradigma}

Al inicio de la entrevista, 2026 se define como el segundo acto de la guerra de los grandes modelos. El primer acto es Chat: el usuario y el modelo interactúan mediante preguntas y respuestas, y la competencia gira en torno a la escala del preentrenamiento, la alineación con instrucciones y la experiencia con contextos cortos. El segundo acto es Agent: el modelo entra en un entorno, invoca herramientas y ejecuta tareas de largo alcance, y la competencia se desplaza hacia el postentrenamiento, la RL infra, la retroalimentación del entorno, el costo y la confiabilidad.

\begin{figure}[H]
\centering
\includegraphics[width=0.94\textwidth]{chat-to-agent-paradigm.png}
\caption{Cambio de paradigma: de un predominio de Chat/Pre-train a un predominio de Agent/Post-train.}
\end{figure}

\begin{knowledgebox}{Lectura de la figura: cómo interpretar esta imagen}
A la izquierda, la etapa Chat pone el acento en las preguntas y respuestas, la recuperación de información y la generación, y el peso de los recursos recae en los datos de preentrenamiento y el modelo base; a la derecha, la etapa Agent pone el acento en las herramientas, el entorno y las tareas de largo alcance, y el peso de los recursos se traslada al rollout, la RL infra, la evaluación y los marcos. La flecha de la figura no dice que el preentrenamiento haya dejado de importar, sino que el foco de la competencia se ha desplazado.
\end{knowledgebox}

\subsection{El papel detonante de OpenClaw}

Al principio, Luo Fuli (罗福莉) veía OpenClaw como Claude Code con una interfaz de usuario añadida y rechazaba por instinto su envoltorio de mercadotecnia. Pero, tras usarlo en serio, lo definió como un marco de Agent que marca época, porque permite que el marco compense las carencias del modelo y que los miembros comunes de un equipo participen, a través del marco, en «elevar el nivel de inteligencia». La importancia de OpenClaw no está en una interfaz concreta, sino en que convierte el marco de Agent en un paradigma de investigación.

\begin{importantbox}{La esencia de OpenClaw}
OpenClaw organiza el modelo, las herramientas, el entorno de la tarea y la imaginación humana dentro de un marco iterable. No es solo una «innovación en la forma del producto», sino que hace explícita la pregunta de investigación de la era del postentrenamiento: cómo lograr que el marco de Agent y el modelo se potencien mutuamente.
\end{importantbox}

\subsection{Por qué 2026 es el año de la transformación de la productividad}

Luo Fuli considera que los desarrolladores chinos reaccionan con más fuerza ante herramientas del tipo de OpenClaw, por un lado porque su necesidad de mejorar la eficiencia es más apremiante, y por otro porque en China hay una gran cantidad de modelos baratos y fáciles de usar. Si una tarea compleja de Agent cuesta 10 yuanes de API pero sustituye un trabajo humano valorado en 1,000 yuanes, el incentivo para adoptarla es muy fuerte. Las condiciones previas de la revolución de productividad de los Agent son un marco lo bastante bueno, modelos lo bastante baratos y un retorno de la tarea lo bastante claro.

\begin{warningbox}{No confundir popularidad con madurez}
Que un marco se vuelva popular puede indicar que respondió a una necesidad de productividad, pero no significa que ya tenga madurez de nivel industrial. La madurez real exige cumplir a la vez con la tasa de finalización de extremo a extremo, la eficiencia de costos, la velocidad, la confiabilidad y los límites de seguridad.
\end{warningbox}

\subsection{Resumen de la sección}

La etapa Chat muestra la inteligencia del modelo; la etapa Agent exige que el modelo entre en el entorno de la tarea. El impacto de OpenClaw consiste en que muchas personas percibieron por primera vez que el marco, las herramientas y las formas de uso colectivo pueden cambiar con rapidez el límite superior visible de la capacidad del modelo.
\section{Postentrenamiento y RL Infra: el sistema central en la era de los Agent}

La entrevista menciona varias veces que el paradigma Agent depende mucho del postentrenamiento. La razón es que un Agent ya no se limita a responder una sola pregunta, sino que realiza acciones de varios pasos dentro de un entorno. Necesita rollout para generar trayectorias, reward/verifier para dar retroalimentación, policy update para modificar la política, evaluation para medir los resultados de extremo a extremo, y además scaffold para gestionar las herramientas y las tareas.

\begin{figure}[H]
\centering
\includegraphics[width=0.94\textwidth]{agent-rl-infra-loop.png}
\caption{El ciclo cerrado de Agent post-train y RL infra.}
\end{figure}

\begin{knowledgebox}{Lectura de la figura: el postentrenamiento de un Agent no es un módulo aislado}
En la figura, el entorno de tareas, el rollout engine, reward/verifier, policy update, eval/cost y scaffold forman un ciclo cerrado. Esto muestra que el postentrenamiento no consiste en «aplicar otra ronda de fine-tuning», sino que es un sistema de ingeniería: cómo se generan las tareas, cómo se recolectan las trayectorias, cómo se define la retroalimentación, cómo se actualiza la política y cómo se miden los costos; todo debe resolverse al mismo tiempo.
\end{knowledgebox}

\subsection{Del motor de inferencia de rollout a un sistema centrado en el Agent}

En la descripción hay una frase decisiva: el sistema pasa de estar «centrado en el motor de inferencia de rollout» a ser un sistema complejo «centrado en el Agent». El primero gira más bien en torno a que el modelo genere respuestas y trayectorias de razonamiento; el segundo requiere el ciclo cerrado completo del entorno. También cambian las exigencias para el equipo: debe ser lo bastante ágil para desarrollar con rapidez una RL infra adaptada al nuevo paradigma.

\noindent\begin{tabularx}{\textwidth}{p{0.22\textwidth}p{0.34\textwidth}X}
\toprule
Etapa & Objeto de entrenamiento/evaluación & Capacidades requeridas del equipo \\
\midrule
Postentrenamiento de Chat & Calidad de las respuestas, preferencias, cadena de razonamiento, alineación & Limpieza de datos, anotación de preferencias, optimización por instrucciones/preferencias. \\
Postentrenamiento de Reasoning & Matemáticas verificables, código, trayectorias de razonamiento & verifier, muestreo, búsqueda, diseño de reward. \\
Postentrenamiento de Agent & Tareas de entorno de varios pasos, cadenas de herramientas, tasa de finalización de extremo a extremo & rollout infra, entornos de tareas, orquestación de herramientas, evaluación de costo y seguridad. \\
\bottomrule
\end{tabularx}

\subsection{Las dificultades del RL scaling}

Cómo hacer bien el RL scaling en los Agent es, para Luo Fuli (罗福莉), una dirección clara en la actualidad, pero que todavía requiere exploración. La dificultad es que el reward de las tareas de un Agent a menudo no es tan claro como las pruebas de código o las respuestas matemáticas; los entornos de tareas son más largos, más abiertos y más propensos a tener estados ocultos; y el costo también es mayor. El RL scaling no consiste solo en ampliar el muestreo, sino en lograr que las tareas, los entornos, la retroalimentación, el costo y la seguridad puedan escalarse.

\begin{importantbox}{Restricciones clave del postentrenamiento de Agent}
El cuello de botella del postentrenamiento de Agent no es simplemente la falta de algoritmos, sino la falta de entornos de tareas y mecanismos de retroalimentación escalables. Sin un entorno estable, el rollout no puede escalarse; sin un verifier confiable, el RL aprende de forma sesgada; sin control de costos, el sistema no puede llevarse a producción.
\end{importantbox}

\subsection{Resumen de la sección}

En la era de los Agent, el postentrenamiento es ingeniería de sistemas. Exige que el equipo pase de la perspectiva del entrenamiento de modelos a la perspectiva del ciclo cerrado del entorno, e integre rollout, reward, policy, evaluation, scaffold y el control de costos en una infraestructura escalable.
\section{MiMo-V2: orquestación de múltiples modelos y emboscada al ecosistema}

Luo Fuli (罗福莉) describe la serie MiMo-V2 como un «despertar» y una «emboscada». Hace un año, cuando trabajaban en multimodalidad y voz, estos modelos funcionaban como modelos independientes; al ver OpenClaw, ella se dio cuenta de pronto de que podían organizarse y orquestarse dentro de un marco de Agent: el modelo de lenguaje se encarga de la planificación y el razonamiento, el modelo multimodal de la percepción, el modelo de voz de la interacción y los modelos pequeños de las llamadas frecuentes de bajo costo.

\begin{figure}[H]
\centering
\includegraphics[width=0.92\textwidth]{mimo-orchestration.png}
\caption{Orquestación de múltiples modelos del ecosistema MiMo-V2 en tareas de Agent.}
\end{figure}

\begin{knowledgebox}{Lectura de la figura: por qué no se reúnen todas las capacidades en un solo modelo}
En la figura, el Agent scaffold ocupa el centro y lo rodean los modelos de lenguaje, multimodal, de voz y pequeños. La conclusión que respalda es que la productividad de extremo a extremo no proviene necesariamente de un único modelo de máximo tamaño, sino de asignar distintos modelos según la etapa de la tarea, para obtener una mejor combinación de velocidad, costo y capacidad.
\end{knowledgebox}

\subsection{Costo, velocidad y precio}

Luo Fuli insiste en varias ocasiones en que una revolución de la productividad debe atender la tasa de finalización de extremo a extremo y la eficiencia de costos. La generación de voz no necesita un gran modelo integrado, y también es dudoso que la comprensión multimodal justifique un modelo más grande. El marco de Agent da más margen de acción a modelos pequeños que no son de primer nivel, pero sí más baratos y rápidos, porque el marco puede compensar sus carencias y estabilizar la salida.

\begin{warningbox}{No preguntes solo «cuál modelo es el más fuerte»}
En un sistema de Agent conviene preguntar qué modelo es el más adecuado para cada etapa. La planificación, la percepción, la voz, la verificación, la ejecución y el resumen pueden requerir modelos distintos. Si el modelo más fuerte es demasiado lento y caro, no necesariamente es la solución óptima para un sistema en producción.
\end{warningbox}

\subsection{Colaboración entre modelos del mismo ecosistema}

Luo Fuli menciona que los modelos entrenados dentro de un mismo ecosistema comparten conocimiento de fondo y límites de capacidad, por lo que se les puede repartir el trabajo en un marco de Agent con más confianza. Lo esencial aquí no es la sinergia de marca, sino que el sistema sepa en qué destaca cada modelo, cuánto cuesta y cómo se conecta su salida con el paso siguiente.

\begin{importantbox}{Implicaciones de ingeniería de MiMo-V2}
En la entrevista, la relevancia de MiMo-V2 no se limita a la publicación de varios modelos: muestra que, en la era de los Agent, una familia de modelos debe orquestarse como un sistema. Los modelos, las herramientas, las tareas y el enrutamiento por costo constituyen en conjunto la capacidad del producto.
\end{importantbox}

\subsection{Resumen de la sección}

La sección sobre MiMo-V2 muestra que, en la era de los Agent, los ecosistemas de múltiples modelos pasarán de la «publicación en paralelo» a la «orquestación de tareas». La competencia futura no consistirá solo en quién tiene el modelo más fuerte, sino también en quién logra organizar varios modelos en un sistema de Agent de bajo costo y alta tasa de finalización.
\section{Asignación de cómputo: el modelo de 1T como boleto de entrada y la proporción 3:1:1}

Uno de los juicios más concretos sobre recursos en la entrevista es la idea de que «un modelo de 1T es el boleto de entrada», junto con el cambio en la proporción de cómputo. Luo Fuli (罗福莉) considera que, si se quiere alcanzar un nivel de Agent cercano al de Claude Opus 4.6, un modelo base de escala 1T puede ser el boleto de entrada. Sin embargo, al entrar en la etapa de Agent, la asignación de cómputo ya no puede quedar acaparada por el preentrenamiento.

\begin{figure}[H]
\centering
\includegraphics[width=0.92\textwidth]{compute-allocation.png}
\caption{Tendencia de la asignación de cómputo desde la etapa de Chat hasta la etapa de Agent.}
\end{figure}

\begin{knowledgebox}{Lectura de la figura: qué significan 3:5:1 y 3:1:1}
Las proporciones 3:5:1 y 3:1:1 de la figura son cifras mencionadas de forma oral en la entrevista y sirven para expresar el desplazamiento del peso de los recursos. En la era de Chat, el preentrenamiento tenía una participación mayor; en la era de Agent, la relación entre research, pre-train y post-train se acerca más al equilibrio, y la proporción de GPU destinada al posentrenamiento aumenta de forma notable.
\end{knowledgebox}

\subsection{En qué sentido 1T es un boleto de entrada}

«Un modelo de 1T es el boleto de entrada» no significa que basten 1T de parámetros, ni que los modelos pequeños carezcan de valor. Apunta a un umbral de capacidad para un Agent de frontera: si se pretende competir con los sistemas de Agent más fuertes, el modelo base debe alcanzar un nivel suficientemente alto en escala, conocimiento, razonamiento, código y capacidad para tareas de varios pasos. El posentrenamiento y el framework pueden amplificar el modelo, pero no pueden crear de la nada todas las capacidades.

\begin{warningbox}{La escala no es la única respuesta}
El número de parámetros es solo una variable de entrada. Un sistema de Agent también depende de los datos de entrenamiento, el contexto, la capacidad de programación, el uso de herramientas, el posentrenamiento, el entorno de la tarea, el costo y la ejecución organizacional. Interpretar «1T como boleto de entrada» como «basta con acumular parámetros para ir a la cabeza» es una lectura errónea.
\end{warningbox}

\subsection{Equilibrio entre post-train y pre-train}

Luo Fuli considera que, en los equipos de primer nivel, la proporción de recursos entre preentrenamiento y posentrenamiento ya se acerca a 1:1. Esto significa que el posentrenamiento ya no es el cierre que sigue al preentrenamiento, sino un frente principal al mismo nivel que este. En las tareas de Agent en particular, el posentrenamiento determina directamente si el modelo puede adaptarse al entorno, a las herramientas y a las tareas de largo alcance.

\begin{importantbox}{La asignación de recursos es una decisión estratégica}
Cómo repartir las GPU no es una cuestión financiera, sino una decisión sobre la ruta técnica. Seguir concentrando todas las GPU en el preentrenamiento, o reservar suficientes para research, post-train, rollout y eval, produce capacidades organizacionales completamente distintas.
\end{importantbox}

\subsection{Resumen de la sección}

En la era de Agent, la asignación de cómputo es más equilibrada y más compleja. La competencia de frontera requiere tanto un modelo base grande como recursos para el posentrenamiento y la investigación. El mensaje central de la proporción 3:1:1 es que el posentrenamiento pasó de ser un actor secundario a ser el protagonista.
\section{Igualdad organizacional: por qué la investigación sobre Agent necesita diversidad}

En la entrevista, la «igualdad organizacional» no es una consigna política, sino un método para organizar la investigación. Luo Fuli (罗福莉) considera que el posentrenamiento necesita diversity, y que la participación de quienes trabajan en preentrenamiento es un buen complemento. Un sistema de Agent abarca el modelo, los datos, el RL, el producto, la ingeniería, el costo, la seguridad y el entorno del usuario; para un equipo de un solo perfil es difícil cubrir todos estos problemas.

\begin{figure}[H]
\centering
\includegraphics[width=0.92\textwidth]{org-equality.png}
\caption{Igualdad organizacional en la investigación sobre Agent: varias funciones entran en el mismo ciclo cerrado de investigación.}
\end{figure}

\begin{knowledgebox}{Lectura de la figura: la igualdad organizacional no es igualitarismo}
En la figura, pre-train, post-train/RL, product/UX, infra/systems y leadership se conectan con Agent research. Esto indica que las distintas funciones deben entrar en el ciclo cerrado de investigación, en lugar de que un único grupo monopolice el juicio. La igualdad significa que la información y la responsabilidad entran en un mismo sistema, no que todos hagan el mismo trabajo.
\end{knowledgebox}

\subsection{Cómo la inteligencia colectiva mejora el framework de Agent}

Durante el Año Nuevo chino, Luo Fuli impulsó que todo el equipo usara OpenClaw. Ella subraya que la imaginación de una sola persona es limitada; cuando los miembros del equipo muestran en el grupo general «resulta que también se puede usar así», despiertan la imaginación de los demás. La inteligencia colectiva no es un concepto abstracto: acelera la investigación al compartir experiencias de uso, estimular mutuamente la imaginación de escenarios y descubrir con rapidez las capacidades y las carencias del framework.

\begin{importantbox}{La organización también es un entorno de entrenamiento}
Si un Agent puede aprender de su entorno, un equipo también puede aprender de un entorno de uso compartido. Hacer que el equipo entre en conjunto a nuevas herramientas, nuevos paradigmas y nuevos entornos de tareas equivale, en esencia, a aplicarle un post-train a la organización.
\end{importantbox}

\subsection{Cultura y valores}

La descripción del video subraya que, en este cambio tecnológico profundo, el fundamento de un AI Lab son su cultura y sus valores. Más adelante, Luo Fuli también comenta que lo que hace cada día debe mejorar un poco el mundo: que el trabajo tedioso se sustituya y que las personas tengan tiempo para hacer cosas de mayor valor. Estos valores influyen en cómo el equipo elige sus tareas, cómo define sus benchmarks y cómo asigna sus recursos.

\begin{warningbox}{Los valores no son un adorno}
A medida que los modelos pueden sustituir cada vez más el trabajo humano, lo que un equipo decide optimizar, sustituir o conservar se vuelve parte de las decisiones técnicas. Sin valores que lo acoten, cuanto más capaz sea un Agent, mayor puede ser también el riesgo.
\end{warningbox}

\subsection{Resumen de la sección}

En la era de los Agent, la capacidad organizacional no consiste solo en contratar a más personas, sino en hacer que personas con distintos perfiles entren juntas en el ciclo cerrado de investigación. La igualdad organizacional, la inteligencia colectiva y los valores determinan si un laboratorio puede adoptar con rapidez un nuevo paradigma.
\section{Arquitecturas alternativas, más allá de 1T e historia de la evolución de los modelos}

La segunda mitad de la entrevista repasa la evolución de los modelos en los últimos tres años: ChatGPT abrió la interacción conversacional de contexto corto; en 2023 los equipos de código abierto se pusieron al día en preentrenamiento; en la etapa de reasoning, el código y las matemáticas se volvieron señales verificables; en la etapa de Agent, el código y la retroalimentación del entorno se convirtieron en la vía de entrada para entrenar tareas de largo alcance. Luo Fuli (罗福莉) considera que la brecha generacional de los equipos chinos en pre-train ya es muy pequeña y que la próxima competencia está en el post-train de Agent y en el RL scaling.

\begin{figure}[H]
\centering
\includegraphics[width=0.96\textwidth]{model-evolution-timeline.png}
\caption{Evolución de los modelos en los últimos tres años: de ChatGPT al postentrenamiento de Agent.}
\end{figure}

\begin{knowledgebox}{Lectura de la figura: cómo cada etapa hereda de la anterior}
La línea de tiempo no es una relación de sustitución. ChatGPT permitió a los usuarios percibir la inteligencia del preentrenamiento; la recuperación del código abierto completó los datos y la arquitectura; reasoning hizo de las tareas verificables una señal de entrenamiento; los prototipos de Agent incorporaron el código y el entorno a las tareas de largo alcance; y 2026 llevó el post-train scaling a la línea principal.
\end{knowledgebox}

\subsection{El sentido de las arquitecturas alternativas}

Luo Fuli menciona decisiones «a contracorriente» como DeepSeek-V2, MoE y las modificaciones de attention. Una arquitectura alternativa no va contra la corriente por el simple hecho de hacerlo, sino que busca investigación escalable en entornos con recursos limitados o sensibles al costo. Si una investigación es solo un paper trick, no necesariamente se convierte en un modelo de nivel industrial; pero si logra escalar y termina por alcanzar un modelo de alto nivel, es frontier research valiosa.

\noindent\begin{tabularx}{\textwidth}{p{0.22\textwidth}p{0.34\textwidth}X}
\toprule
Ruta & Problema que resuelve & Criterio de industrialización \\
\midrule
MoE & Reducir el costo de cómputo por token mediante activación dispersa & Requiere que el enrutamiento, el balanceo de carga, la estabilidad del entrenamiento y la eficiencia del despliegue se cumplan a la vez. \\
Modificación de attention & Reducir el costo de contexto largo y de inferencia & Hay que ver si mantiene la calidad en tareas reales, no solo en un benchmark local. \\
Modelo pequeño + marco & Usar un modelo barato junto con un Agent scaffold & Adecuado para etapas de alta frecuencia, bajo costo y cuyas carencias puede suplir el marco. \\
Base de 1T & Proporcionar la base de capacidades de Agent de frontera & Es el boleto de entrada, no la respuesta completa. \\
\bottomrule
\end{tabularx}

\subsection{La IA no tiene crisis de supervivencia}

Luo Fuli compara la evolución humana con la evolución de los modelos: los seres humanos evolucionaron en un entorno natural y bajo presión de supervivencia; los modelos no tienen una crisis de supervivencia propia, pero cuentan con mucho cómputo, con el conocimiento humano y con la ayuda de las personas. Por ello, la trayectoria evolutiva de la IA puede ser más libre, más dispersa, más rápida y también más distinta de la inteligencia biológica. Este pasaje nos recuerda que no debemos aplicar de forma forzada la analogía de la evolución humana a los modelos.

\begin{knowledgebox}{Los límites de la analogía}
«El lenguaje apareció muy tarde en la inteligencia humana» puede ayudarnos a entender la estructura de triángulo invertido de la IA, pero el entorno de los modelos es completamente distinto del humano. El siguiente paso de la IA no necesariamente reproducirá la trayectoria del cuerpo o de los sentidos humanos; es posible que primero siga expandiéndose en coding, en productividad y en entornos digitales.
\end{knowledgebox}

\subsection{Negar cada día al yo de ayer}

Luo Fuli dice que en el último medio año casi todos los días ha negado a su yo del día anterior. Esta expresión describe bien la etapa de Agent de 2026: las herramientas, los marcos, los modelos, los costos y las formas de organización cambian con rapidez. Para un equipo de investigación, la estabilidad no consiste en aferrarse a juicios antiguos, sino en construir mecanismos que permitan actualizarlos con rapidez.

\begin{importantbox}{De dónde viene la velocidad de investigación}
La velocidad de investigación no proviene solo de personas inteligentes, sino también del entorno: cómputo suficiente, infraestructura suficientemente buena, retroalimentación suficientemente rápida, una organización suficientemente abierta y valores suficientemente claros. Que el entorno importe más que la experiencia es el juicio central con el que cierra este episodio.
\end{importantbox}

\subsection{Resumen de la sección}

En los últimos tres años, la evolución de los modelos pasó de la interacción conversacional, la recuperación del código abierto y la verificación mediante reasoning al postentrenamiento de Agent. El liderazgo en la siguiente etapa dependerá de arquitecturas que puedan escalar, de los sistemas de postentrenamiento, de la asignación de recursos y de la capacidad de negarse a sí mismo con rapidez.
\section{Consenso y competencia actuales: cómo se ponen al día los equipos chinos}

Luo Fuli (罗福莉) considera que ya existe consenso en que el camino de Anthropic era el correcto; una vez que el camino quedó más claro, los equipos chinos de modelos grandes entraron en una fase de recuperación acelerada. Hoy la brecha generacional en pre-train es muy pequeña, e incluso en algunas arquitecturas los equipos chinos tienen ventaja; donde de verdad hay que ir all in es en el Agent post-train, en particular en cómo hacer bien el RL scaling sobre Agents.

\subsection{Por qué Coding importa en cada paradigma}

Coding ha tocado puntos decisivos en varias etapas. En la etapa de preentrenamiento, los datos de código tienen alta calidad y una estructura clara; en la etapa de reasoning, el código y las matemáticas ofrecen retroalimentación verificable; en la etapa de Agent, el desarrollo de software es en sí mismo una tarea de largo horizonte, con entornos reales, pruebas, depuración y dependencias complejas. Es a la vez señal de entrenamiento y escenario de producto.

\begin{importantbox}{El triple papel de Coding}
Coding es datos, evaluación y producto. Como datos, es estructurado y de alta calidad; como evaluación, cuenta con pruebas y retroalimentación de ejecución; como producto, corresponde directamente a escenarios de productividad de alto valor. Por eso Coding importa en cada paradigma.
\end{importantbox}

\subsection{El entorno importa más que la experiencia}

En la parte final, Luo Fuli enfatiza que el entorno importa más que la experiencia. Por entorno entiende la capacidad de cómputo, la infraestructura, los chips, el sistema operativo, el tráfico, las redes sociales, la cultura organizacional y los recursos estratégicos. Que un equipo pueda liderar en el nuevo paradigma depende de si logra que todos estos recursos se adapten en conjunto a la arquitectura de Agent.

\begin{warningbox}{La experiencia puede convertirse en una carga}
Cuando el paradigma cambia con rapidez, la experiencia previa ayuda a juzgar, pero también limita la imaginación. Lo que de verdad importa es si el entorno le da al equipo retroalimentación y aceleración continuas, de modo que pueda refutar una y otra vez sus juicios anteriores y reconstruir sus procesos anteriores.
\end{warningbox}

\subsection{Resumen de la sección}

El consenso de la competencia actual es: la brecha en pre-train se reduce y el Agent post-train se vuelve la línea principal. La oportunidad de los equipos chinos está en que el camino es más claro, son sensibles al costo y la retroalimentación de las aplicaciones es rápida; los retos están en la RL infra, los entornos de tareas y la reestructuración organizacional.
\section{Apéndice: índice de términos y ruta de repaso}

\subsection{Términos clave: índice de palabras clave de este episodio}

\noindent\begin{longtable}{p{0.23\textwidth}p{0.32\textwidth}p{0.36\textwidth}}
\toprule
Término & Explicación en una frase & Papel en este episodio \\
\midrule
OpenClaw & Producto y framework de código abierto que detonó el consenso sobre los frameworks de Agent & Representa la exhibición del paradigma de Agent convertido en producto. \\
Post-train & Entrenamiento que continúa después del preentrenamiento y gira en torno a tareas, preferencias, entornos y retroalimentación & Principal terreno de competencia en la era de los Agent. \\
RL Infra & Infraestructura que sostiene el rollout, el reward, el policy update y la eval & Determina si el posentrenamiento de los Agent puede escalar (scale). \\
Rollout & Generación de trayectorias y procesos de interacción en el entorno de la tarea & Fuente de los datos y de la retroalimentación de RL. \\
MiMo-V2 & Familia de modelos de Xiaomi (小米), caso central de la entrevista & Muestra cómo un framework de Agent orquesta varios modelos. \\
Boleto de entrada de 1T & Umbral de escala del modelo base necesario para equipararse con las capacidades de Agent de frontera & Explica por qué el modelo base sigue siendo importante. \\
Igualdad organizacional & Varias funciones se integran en un mismo ciclo de investigación & Explica por qué el posentrenamiento necesita diversity. \\
El entorno importa más que la experiencia & La velocidad de investigación proviene de la capacidad de cómputo, la infra, la organización y el entorno de retroalimentación & Juicio estratégico del cierre. \\
\bottomrule
\end{longtable}

\subsection{Ruta de repaso}

Se recomienda leer primero el capítulo 2 para entender la transición de paradigma de Chat a Agent; después, el capítulo 3 para dominar el ciclo de posentrenamiento, y luego los capítulos 4--6 para entender la orquestación de MiMo, la asignación de capacidad de cómputo y la igualdad organizacional. Los capítulos 7--8 sirven como juicio sobre la industria: por qué en 2026 la competencia ya no se limita al modelo base, sino que consiste en la combinación de arquitectura, posentrenamiento, entorno y capacidad organizacional.

\begin{knowledgebox}{Conexión con EP139}
La entrevista con Su Yu (苏煜) en EP139 ofrece la historia técnica de los Agent y el framework de universal digital agent; la entrevista con Luo Fuli (罗福莉) en EP138 muestra cómo un laboratorio de modelos convierte ese framework en posentrenamiento, orquestación de modelos, asignación de recursos y reestructuración organizacional. Ambos episodios deben leerse juntos.
\end{knowledgebox}

\subsection{Resumen de la sección}

Las palabras clave de este episodio pueden condensarse en seis: OpenClaw, posentrenamiento, RL infra, orquestación de varios modelos, reasignación de capacidad de cómputo e igualdad organizacional. Todas apuntan a un mismo juicio: en la era de los Agent, la competencia es una competencia entre sistemas.

\section{Detalles de entrenamiento, costos y arquitecturas alternativas: por qué esto es ingeniería de sistemas}

Aunque la sección «Detalles de entrenamiento y costos» de la entrevista es breve, es la clave para entender todo el episodio. Luo Fuli (罗福莉) insiste en que, en la era de los Agent, todos los caminos deben volver a la tasa de finalización de extremo a extremo, la velocidad y el costo. Que un modelo sea más fuerte en un benchmark no significa que funcione mejor dentro de un framework de Agent; que un modelo sea más grande tampoco significa que valga la pena invocarlo en todas las etapas. La verdadera ingeniería de sistemas consiste en reunir en una misma tabla la capacidad del modelo, las etapas de la tarea, las llamadas a herramientas, la latencia, el precio y la confiabilidad para tomar decisiones.

\begin{importantbox}{El costo no es una partida financiera secundaria, sino parte de la capacidad del modelo}
Si una tarea de Agent requiere decenas de rondas de llamadas al modelo, ejecución de herramientas y retroalimentación del entorno, el costo de cada llamada se multiplica a lo largo de la tarea de horizonte largo. Los modelos pequeños y baratos, los modelos especializados, el enrutamiento entre varios modelos, la caché, los verificadores y la orquestación del framework influyen directamente en que el sistema pueda ponerse en producción.
\end{importantbox}

\subsection{Términos clave: vocabulario de entrenamiento y costos}

\noindent\begin{tabularx}{\textwidth}{p{0.23\textwidth}p{0.31\textwidth}X}
\toprule
Término & Problema que resuelve & Relación con este episodio \\
\midrule
Tasa de finalización de extremo a extremo & Si la tarea completa termina cumpliéndose & Un producto de Agent no puede evaluarse solo por acciones parciales o respuestas de una sola ronda. \\
Enrutamiento de llamadas & Elegir modelos distintos para etapas distintas & Usar un modelo grande para el razonamiento crítico y uno pequeño para las etapas frecuentes y de bajo riesgo. \\
Verificador & Determinar si el resultado de la tarea es correcto & Es esencial para el RL feedback y la evaluación de confiabilidad. \\
Costo de Rollout & Costo de cómputo y de API de generar múltiples trayectorias & Determina si el postentrenamiento puede hacer scale. \\
Latency & Tiempo de espera del usuario y duración total de la tarea & Influye en que una herramienta de productividad sea utilizable. \\
Scaffold & Framework externo del Agent, orquestación de herramientas y gestión de estado & Puede compensar las deficiencias del modelo y también elevar el techo de los modelos de punta. \\
\bottomrule
\end{tabularx}

\subsection{Por qué importan las arquitecturas alternativas}

Aquí las arquitecturas alternativas no son «trucos ingeniosos sin sustancia», sino investigación escalable bajo restricciones de recursos. DeepSeek-V2, MoE, las modificaciones de attention y otras líneas de trabajo plantean, en esencia, la misma pregunta: cómo obtener más inteligencia por unidad de costo con un cómputo y un hardware dados. En la era de los Agent esta pregunta se vuelve más apremiante, porque las tareas de horizonte largo invocan el modelo una y otra vez, y cualquier diferencia de eficiencia en un solo paso se amplifica.

\begin{warningbox}{No hay que ver las arquitecturas alternativas solo como innovación de artículos}
Si una arquitectura no puede hacer scale hasta un modelo industrial, no puede desplegarse en un sistema de inferencia real o no logra reducir el costo de extremo a extremo de las tareas, su valor para los productos de Agent es limitado. El «nivel de calidad industrial» que subraya la entrevista exige justamente que la innovación arquitectónica termine incorporándose a sistemas que funcionen.
\end{warningbox}

\subsection{Resumen de la sección}

Los detalles de entrenamiento, los costos y las arquitecturas alternativas muestran en conjunto que la ruta tecnológica de la era de los Agent no es una competencia de una sola variable por «modelos más grandes», sino una competencia por la eficiencia del sistema de extremo a extremo. La estructura del modelo, el postentrenamiento, el framework y el costo de inferencia deben optimizarse en conjunto.
\section{Mentalidad de investigación: negar cada día a quien uno fue ayer}

Luo Fuli (罗福莉) dice que lo que más ha sentido en el último medio año es que «cada día niega a quien fue ayer». No es una expresión emocional, sino una mentalidad de investigación ante un cambio de paradigma acelerado: cuando OpenClaw, Claude Code, el posentrenamiento de Agent y la orquestación de múltiples modelos aparecen con rapidez, los juicios anteriores pierden validez en cuestión de semanas o incluso de días. El equipo debe ser capaz de actualizar hipótesis, rehacer experimentos y reajustar la asignación de recursos con rapidez.

\subsection{El cambio de reward: de la inversión cuantitativa a los modelos grandes}

En la entrevista se menciona que, en la inversión cuantitativa, el precio es el reward, que resulta relativamente claro; en cambio, en el campo de los modelos grandes el reward es menos claro y más cambiante. Con los Agent esto se acentúa: el éxito de la tarea, el valor para el usuario, los límites de seguridad, la eficiencia en costos y el aprendizaje a largo plazo pueden formar parte del reward. Esto exige que el investigador no se limite a optimizar un solo número, sino que defina de manera continua «qué es lo bueno».

\begin{knowledgebox}{Cómo investigar cuando el reward no es claro}
Cuando el reward no es claro, la investigación no puede detenerse; primero hay que construir indicadores sustitutos verificables: tasa de finalización de tareas, preferencia humana, costo, tipos de falla, capacidad de recuperación, incidentes de seguridad, retención de usuarios, entre otros. Después hay que comprobar de forma constante si esos indicadores sustitutos realmente sirven al valor final.
\end{knowledgebox}

\subsection{El entorno importa más que la experiencia}

«El entorno importa más que la experiencia» es el juicio estratégico con el que cierra este episodio. Aquí el entorno incluye la capacidad de cómputo, la infraestructura, los chips, los sistemas operativos, el tráfico, la dimensión social, la cultura organizacional y los recursos estratégicos. La experiencia es útil, pero si el entorno no da retroalimentación, no permite equivocarse y probar, no proporciona capacidad de cómputo ni apoya la colaboración entre equipos, la experiencia se convierte pronto en un lastre del paradigma anterior.

\begin{importantbox}{El entorno es la aceleración de la investigación}
La experiencia determina la posición inicial; el entorno determina la aceleración. En la era de los Agent, los cambios son demasiado rápidos: quien obtenga antes retroalimentación real, más rollout, mejor infra y una organización más abierta tendrá más probabilidades de liderar en el nuevo paradigma.
\end{importantbox}

\subsection{Resumen de la sección}

Negar cada día a quien uno fue ayer no significa carecer de juicio, sino construir un sistema capaz de actualizar sus juicios. La mentalidad de investigación en la era de los Agent consiste en enlazar valores, retroalimentación experimental, entorno de recursos y ajustes organizacionales en un ciclo continuo de aprendizaje.
\section{Comparación con EP139: del marco de historia técnica a la acción del laboratorio}

La entrevista con Su Yu (苏煜) en EP139 ofrece una historia técnica de los Agent: cómo el Agent, entendido como un sistema con entidad, entorno y acción orientada a objetivos, pasó de logical agent y semantic parsing a language agent y universal digital agent. La entrevista con Luo Fuli (罗福莉) en EP138 responde otra pregunta: cuando este paradigma llega al interior de un laboratorio de modelos, cómo debe actuar el equipo.

\noindent\begin{tabularx}{\textwidth}{p{0.22\textwidth}p{0.34\textwidth}X}
\toprule
Dimensión & EP139 Su Yu & EP138 Luo Fuli \\
\midrule
Nivel de atención & Genealogía técnica y marco conceptual de los Agent & Recursos, organización y línea de modelos de un AI Lab. \\
Giro central & El OpenClaw Moment forma el consenso sobre los Agent & OpenClaw cambia el paradigma de investigación del equipo. \\
Mecanismos clave & Language Agent, disolución de fronteras, continual learning & Post-train, RL infra, orquestación de varios modelos, asignación de GPU. \\
Problemas de implementación & reliability, cost, deployment, world model & Tasa de finalización de extremo a extremo, costo y velocidad, igualdad organizativa, retroalimentación del entorno. \\
\bottomrule
\end{tabularx}

\begin{knowledgebox}{Cómo leerlas en conjunto}
Primero lee EP139 para entender por qué el Agent es un problema técnico de largo plazo; después lee EP138 para entender cómo un equipo de modelos ajusta sus recursos, su sistema de entrenamiento y su estructura organizativa una vez que ese problema técnico se convierte en consenso de la industria.
\end{knowledgebox}

\subsection{Resumen de la sección}

Los dos episodios forman un par: EP139 es el mapa conceptual de los Agent y EP138 es el mapa de acción de un laboratorio de modelos en la era de los Agent. La cola de Zhang Xiaojun AI puede seguir organizándose a partir de estos dos mapas.
\section{Apéndice II: lista de acciones de un AI Lab ante el paradigma Agent}

Esta entrevista puede convertirse en una lista de acciones para equipos de modelos. No pretende que cada empresa la copie tal cual, sino ayudar al lector a llevar el «cambio de paradigma» a preguntas que se puedan ejecutar.

\noindent\begin{longtable}{p{0.25\textwidth}p{0.34\textwidth}p{0.32\textwidth}}
\toprule
Decisión & Pregunta que hay que responder & Práctica equivocada \\
\midrule
Marco de Agent & ¿Cómo se conecta nuestro modelo con herramientas, entornos, memoria y tareas de largo horizonte? & Tratar al Agent solo como un system prompt complejo. \\
RL infra & ¿Cómo se escalan rollout, reward, verifier y eval? & Trabajar solo con datos de preferencias fuera de línea, sin construir el ciclo cerrado con el entorno. \\
Recursos de GPU & ¿Cómo se reparten entre research, pre-train y post-train? & Seguir dejando que el preentrenamiento acapare todos los recursos estratégicos. \\
Familia de modelos & ¿Qué etapas usan modelos grandes y cuáles usan modelos pequeños o especializados? & Llamar al modelo más grande para cualquier tarea, sin distinción. \\
Costo y velocidad & ¿Cómo se optimizan en conjunto la tasa de finalización de extremo a extremo, la latencia y el precio? & Mirar solo el benchmark, sin calcular el costo de las tareas reales. \\
Estructura organizacional & ¿Entran juntos al ciclo cerrado el preentrenamiento, el posentrenamiento, el producto y los sistemas? & Dejar que un solo equipo monopolice todas las decisiones. \\
Valores & ¿Qué trabajos se sustituyen y qué valores humanos se conservan? & Optimizar solo métricas destructivas o de corto plazo. \\
\bottomrule
\end{longtable}

\begin{importantbox}{Para qué sirve la lista de acciones}
Si un equipo afirma que va a ir all in Agent, pero no puede responder las preguntas de la tabla anterior, es muy probable que solo esté siguiendo la moda. Una transformación real hacia Agent se refleja en la asignación de recursos, la infraestructura de entrenamiento, el enrutamiento entre modelos, la colaboración organizacional y la elección de valores.
\end{importantbox}

\subsection{Resumen de la sección}

El paradigma Agent no es un lema estratégico, sino un conjunto de decisiones concretas. El lector puede usar esta lista de acciones para comprobar si cualquier laboratorio de modelos, empresa emergente o equipo de producto ha entrado de verdad en la era de los Agent.
\section{Apéndice III: tarjeta de autoevaluación para laboratorios de modelos}

Para convertir esta entrevista en criterios aplicables, puede elaborarse una tarjeta de evaluación simplificada para cualquier laboratorio de modelos o equipo emprendedor de Agent. No es una herramienta de clasificación, sino una ayuda para que el lector juzgue si un equipo comprende de verdad el paradigma Agent.

\noindent\begin{longtable}{p{0.24\textwidth}p{0.33\textwidth}p{0.34\textwidth}}
\toprule
Dimensión & Desempeño de puntaje alto & Señales de puntaje bajo \\
\midrule
Juicio sobre el paradigma & Explica con claridad las diferencias de entrenamiento y de producto entre las tres etapas: Chat, Reasoning y Agent & Concibe Agent solo como llamadas a herramientas o como un prompt complejo. \\
Sistema de posentrenamiento & Cuenta con un ciclo cerrado de entornos de tareas, rollout, verifier, eval y monitoreo de costos & Solo dispone de datos fuera de línea o de ejemplos seleccionados manualmente. \\
Orquestación de modelos & Enruta modelos grandes y pequeños, modelos multimodales y modelos especializados según la etapa de la tarea & Todas las tareas llaman al mismo modelo, el más grande. \\
Asignación de cómputo & research, pre-train y post-train tienen proporciones claras y un ritmo de iteración definido & Los recursos quedan atados a la trayectoria histórica y el posentrenamiento no tiene suficientes GPU. \\
Mecanismos organizacionales & Preentrenamiento, posentrenamiento, producto, sistemas y evaluación participan juntos en las decisiones & Un solo equipo decide a puerta cerrada; la retroalimentación de usuarios reales y de ingeniería no llega. \\
Conciencia de costos & Optimiza en conjunto la tasa de finalización de extremo a extremo, la latencia y el precio & Solo le interesa exhibir capacidades, no la economía por tarea. \\
Valores y límites & Define con claridad qué tareas deben automatizarse y en cuáles las personas deben conservar el juicio & Solo persigue la sustitución y el crecimiento, sin discutir seguridad, responsabilidad ni consecuencias sociales. \\
\bottomrule
\end{longtable}

\begin{knowledgebox}{Cómo usar esta tarjeta de evaluación}
Al leer las siguientes entregas de la serie Zhang Xiaojun AI, pueden situarse las opiniones de cada invitado en estas siete dimensiones. Por ejemplo, la entrevista con Su Yu (苏煜) se inclina más hacia «juicio sobre el paradigma» y «aprendizaje a largo plazo», mientras que la entrevista con Luo Fuli (罗福莉) se inclina más hacia «sistema de posentrenamiento», «orquestación de modelos», «asignación de cómputo» y «mecanismos organizacionales». Así, las notas de varias entregas irán formando un mapa comparable del sistema operativo de un AI Lab.
\end{knowledgebox}

\subsection{Resumen de la sección}

La tarjeta de evaluación convierte los juicios de la entrevista en un marco reutilizable. Recuerda al lector que, en la era de los Agent, no triunfa una sola tecnología aislada: varias dimensiones del sistema deben cumplir el estándar al mismo tiempo.
\section{Síntesis y ampliación}

\subsection{Conclusiones centrales}

\begin{enumerate}
\item La era de los Agent no consiste en una mejora de la capacidad de conversación, sino en que el modelo entra en entornos, herramientas y tareas de largo alcance.
\item El valor de OpenClaw está en que generó un consenso en torno al framework y permitió a los equipos ver que el modelo y el scaffold pueden potenciarse mutuamente.
\item El postentrenamiento, en particular el Agent RL scaling, se convierte en la línea principal, y requiere entornos de tareas, rollout, reward, eval y control de costos.
\item La importancia de MiMo-V2 está en la orquestación de múltiples modelos, no en el lanzamiento de un modelo individual.
\item La asignación de capacidad de cómputo pasa del predominio del preentrenamiento a un reparto más equilibrado entre research, pre-train y post-train.
\item La horizontalidad organizacional, la diversidad y los valores culturales influyen en la rapidez con que un AI Lab adopta un nuevo paradigma.
\item Las oportunidades de los equipos chinos para alcanzar a los líderes están en el Agent post-train, la sensibilidad a los costos y la retroalimentación de las aplicaciones; los retos están en la RL infra y en la agilidad organizacional.
\end{enumerate}

\subsection{Preguntas abiertas}

\begin{itemize}
\item ¿Cuál será la forma predominante de reward y de entorno en el Agent post-train?
\item ¿Seguirá siendo a largo plazo un modelo base de 1T el requisito de entrada, o será desplazado gradualmente por modelos pequeños + framework?
\item ¿La orquestación de múltiples modelos desembocará en un ecosistema unificado o en el enrutamiento de modelos dentro de un mercado abierto?
\item ¿Cómo puede la horizontalidad organizacional evitar convertirse en una colaboración caótica y aumentar de verdad la velocidad de la investigación?
\item Cuando el modelo no tiene una «crisis de supervivencia», ¿cómo deben los seres humanos diseñar su función de valor y su lugar en la sociedad?
\end{itemize}

\subsection{Lecturas adicionales}

\begin{itemize}
\item EP139, Su Yu (苏煜), historia técnica de los Agent: para entender el contexto histórico del OpenClaw Moment y de los Language Agent.
\item Materiales sobre Claude Code / OpenClaw / Computer Use Agent: para entender cómo los frameworks de Agent cambian la forma de los productos.
\item Artículos sobre RLHF, RLAIF, verifier, rollout y agent benchmark: para entender la infraestructura del postentrenamiento.
\item MoE, attention alternatives, multi-modal model orchestration: para entender los problemas de sistemas de modelos detrás de MiMo-V2.
\end{itemize}

\begin{importantbox}{Valoración final}
Quién gane en la era de los Agent no lo decidirá un único modelo aislado. Se parece más a una competencia de sistemas: quien logre integrar más rápido en un ciclo cerrado el modelo, el framework, los entornos de tareas, el postentrenamiento, la capacidad de cómputo, el producto y la organización tendrá más probabilidades de ir a la cabeza en la siguiente etapa.
</br>
\end{importantbox}

\end{document}
